\section{Cube estimates at a fixed positive amplitude}\label{sec:compact}
Throughout this section $a\in[a_0,a_1]$ and $r\in[r_0,r_1]$, with
all four endpoints fixed as in Theorem~\ref{thm:compact}.
Constants depend only on those endpoints unless another dependence is stated.

\subsection{Elementary norm and entropy estimates}
For $|x|\le\gamma<1$ and $y\in[-1,1]$, let
$B_\psi(y,x)=\psi(y)-\psi(x)-V(x)(y-x)$. Then
\begin{equation}\label{eq:bregman}
 \tfrac12(y-x)^2\le B_\psi(y,x)\le U_\gamma(y-x)^2,
 \quad |B_\psi(x+z,x)-\tfrac12\psi''(x)z^2|
 \le C_{\gamma,p}|z|^p\quad(2<p\le3).
\end{equation}
The lower bound is $\psi''\ge1$. The other two follow from interior
Taylor expansion near $z=0$ and continuity on the remaining compact
set, including $y=\pm1$. In particular no endpoint derivative is used.
Put
\[
 p=\min\{3,1+r_1^{-2}\}>2,\qquad
 \theta=(1+r_1^2)^{-1}-\tfrac12>0.
\]
Real hypercontractivity \cite{Bonami,Beckner} gives
\begin{equation}\label{eq:hyper}
 \|T_rz\|_p\le\|z\|_2,\qquad
 \|T_rz\|_2\le P^\theta\|z\|_2
 \quad\text{if }\Pr(z\ne0)\le P.
\end{equation}
For clarity this applies to real functions, not only indicators.
A two-point proof starts with $M(t)=\E|A+tBX|^q$, $1<q\le2$:
the second derivative of $M^{2/q}$ is a nonnegative square plus
$2(q-1)B^2 M^{2/q-1}\E|A+tBX|^{q-2}\ge2(q-1)B^2$.
Jensen proves the last inequality; continuity across its integrable
zeros permits integration. Tensorization and duality give
\eqref{eq:hyper}.
Also, if $z$ has support mass $P$,
\begin{equation}\label{eq:lowproj}
 \|\Pi_{\le1}z\|_2\le C P^{1/4}\|z\|_2.
\end{equation}
Indeed degree-at-most-one polynomials of norm one have uniformly
bounded fourth moments. Duality and H\"older prove the assertion.

\subsection{The semigroup budget and return to the original kernel}
Let $P_t=T_{e^{-t}}$, $t=-\log r$, and set
\[
 y(s)=\cI(P_su),\quad w(s)=\psi(ae^{-s}),\quad z(s)=w(s)-y(s),
 \quad R(s)=D(P_su)-F(y(s)).
\]
For $s>0$, a nonconstant finite-cube kernel becomes strictly interior.
The flow identity gives $z'=R-A(s)z$, where
$2\le A(s)\le K:=\max_{[0,\psi(a_1)]}F'$ is the secant slope of $F$.
Here $R\ge0$, $z\ge0$ by the assumed inequalities. Integration from
positive time and then continuity at zero give
\begin{equation}\label{eq:budget}
 \Delta_r(u)=\int_0^t e^{-\int_s^t A(v)\,dv}R(s)\,ds,
 \quad \int_0^tR(s)\,ds\le e^{Kt}\Delta_r(u),
 \quad z(s)\le e^{Kt}\Delta_r(u).
\end{equation}
No finite production at the original endpoint-valued kernel is required.

For a coordinate oriented so that $b=\E uX_i\ge0$, strong convexity
of $\psi$ at $bX_i$ gives
\[
 \|u-bX_i\|_2^2\le2\{\psi(a)-\psi(b)+\psi(m)\},\qquad m=\E u.
\]
If $\|T_\lambda u-\lambda aX_i\|_2^2=q$, then $m^2\le q$ and
$|b-a|\le\sqrt q/\lambda$. Thus, for $\lambda\ge r_0$ and small $q$,
\begin{equation}\label{eq:return}
 \|u-aX_i\|_2^2
 \le 2V((1+a_1)/2)\sqrt q/r_0+2Lq+q/r_0^2.
\end{equation}
We used the orthogonal identity
$\|u-aX_i\|^2=\|u-bX_i\|^2+(a-b)^2$.
This is an entropy estimate, not an inverse-continuity assertion for noise.

\begin{lemma}[Local stability]\label{lem:local}
There are $\varepsilon_*,c_*>0$ such that
$\|u-saX_i\|_2\le\varepsilon_*$ and $\cI(u)=\psi(a)$ imply
$\Delta_r(u)\ge c_*\|u-saX_i\|_2^2$.
\end{lemma}
\begin{proof}
Write $v=aX_1$, $h=u-v$, $m=\E h$, and $k=rV(ra)/V(a)$.
The exact Bregman expression for
$G=\Delta_r(u)+k(\cI(u)-\psi(a))$ is
\[
 G=k\E B_\psi(v+h,v)-\E B_\psi(rv+T_rh,rv)+(1-k)\psi(m).
\]
Its quadratic form is
\[
 Q(h)=\frac{k\|h\|_2^2}{1-a^2}
 -\frac{\|T_rh\|_2^2}{1-r^2a^2}+(1-k)m^2.
\]
Every nonconstant Walsh coefficient has coefficient at least
$q=k/(1-a^2)-r^2/(1-r^2a^2)$; the constant coefficient is $a^2q$.
The function $(1-t^2)V(t)/t$ has derivative
$[t-(1+t^2)V(t)]/t^2<0$, proving $q>0$. Moreover $0<k<r^2<1$.
Consequently $Q\ge c_H\|h\|^2$, where $c_H=a_0^2\min q>0$.

Choose $\eta_0\le(1-a_1)/2$ with $D_3\eta_0/6\le c_H/8$,
where $D_3$ bounds $|\psi'''|$ on the resulting interior interval.
Split $h=g+b$ on the sets $|h|\le\eta_0$ and $|h|>\eta_0$.
The latter has mass $P\le\|h\|^2/\eta_0^2$.
Taylor expansion on the good set and strong convexity on the other give
\[
 k\E B_\psi(v+h,v)\ge
 \frac{k\|g\|^2}{2(1-a^2)}-\frac{c_H}8\|g\|^2
 +\frac{k_{\min}}2\|b\|^2.
\]
The output remainder is at most $C\|h\|^p$ by
\eqref{eq:bregman}--\eqref{eq:hyper}. Use $\psi(m)\ge m^2/2$.
The good quadratic terms equal $Q(g)/2$. With
$B_*=(1-r_1^2a_1^2)^{-1}$, the bad square costs at most
$B_*P^{2\theta}\|b\|^2/2$ and the cross terms at most
$(B_*P^\theta+\sqrt P)\|g\|\|b\|$.
Young's inequality bounds the latter by
$c_H\|g\|^2/8+2(B_*P^\theta+\sqrt P)^2\|b\|^2/c_H$.
Choose $P_*>0$ so that these two bad coefficients sum to at most
$k_{\min}/4$. Choose $\varepsilon_*$ so that
$\varepsilon_*^2/\eta_0^2\le P_*$ and
$C\varepsilon_*^{p-2}\le\min(c_H,k_{\min})/8$.
The remaining positive coefficients prove the claim. The input values
on the bad set have never been expanded.
\end{proof}

\subsection{Equality and compactness for a joining law}
For the notation in Assumption~\ref{ass:inputs}(5), write
\[
 \mathcal J(\mu)=\E c(A,B)
 +\tfrac12\{\Phi_a(J_a)+\Phi_b(J_b)\}-\Phi_m(J).
\]

\begin{lemma}\label{lem:joining}
For a law $\mu$ of finite cost with $(a-b)/2>0$, $\mathcal J(\mu)=0$
if and only if $\mu$ is a point mass at $(t,-t)$, $0<t<1$.
On the set of laws with $(a-b)/2\ge\beta>0$, $j_0\le J\le j_1<L$,
and $\E c\le C_0$, a joining surplus tending to zero forces
\[
 \tfrac12\E[(A-t)^2+(B+t)^2]\longrightarrow0,
 \qquad t=\psi^{-1}(J),
\]
uniformly over that set.
\end{lemma}
\begin{proof}
Finite cost and positive separation exclude zero child entropy: if
$H(A)=0$ almost surely, then $A$ is an endpoint, and finite cost forces
$B=A$ almost surely, contradicting separation. The same proves positive
entropy for $B$ and excludes endpoint means.
If the parent maximum is its reserve, including a tie, perspective
pooling gives $\mathcal J\ge\E c-W\ge\cK-W$.
Let its optimizer have core mass $k$ and half-difference $d_c$.
Compose \eqref{eq:reserveloss} with retentions
$\tau_1=k/(k+\pi_+)$ and $\tau_2=k+\pi_+$. The loss is at least
\[
 \frac{\delta_0d_c^2}{8}
 \frac{k\pi_++k^2\pi_-}{k+\pi_+}
 \ge\frac{\delta_0}{8}(1-k)d^2.
\]
The deterministic core is paid by the capacity inequality:
$W_c\le(1-m_c^2)F(J_c/(1-m_c^2))\le c_c$.
If $m_c\ne0$, then $Lm_c^2-\psi(m_c)>0$ by the positive entropy
series, and the first inequality is strict: the reserve argument is
strictly below $J_c/(1-m_c^2)>0$, where $S$ strictly increases.
Equality therefore forces $k=1$ and $m_c=0$.
It also forces $\E c=\cK$. Uniqueness of the priced optimizer, at
its exact entropy budget, then forces the actual law to be this point mass.

For a strict $F$ parent, $m\ne0$ and $J<L/2$ by
\eqref{eq:high}. Reflect to $m>0$.
Convexity of $F$ at the original children gives
$\mathcal J\ge\cK-[F(J)-F(j)]$, $j=(J_a+J_b)/2$.
For a one-positive-corner optimizer, \eqref{eq:thinloss} is strict
if $k<1$. If $k=1$, strict convexity of $F$ and the capacity inequality
give $c_c>F(J_c)$ because $m>0$. For a reflected optimizer,
\eqref{eq:reflected} implies $k<1$. Increase $m$ to $1-k$, keeping
$d,e,k,t$ fixed. Its information decreases, stays below $L/2$, and
the demand increases by \eqref{eq:entropyratio}. At that wall the
strict one-corner payment applies. Equality is impossible.
Conversely a point mass $(t,-t)$ has zero child information and cost
$F(\psi(t))$, so its surplus is zero.

Finally, probability laws on the fixed compact square are weakly
compact. The cost sublevel is closed by lower semicontinuity; means
and entropy moments are continuous. The constraints in the lemma thus
define a compact set. The preceding endpoint exclusion applies at every
point of it, so compactness supplies a common positive child entropy
and interior mean margin. All profile terms are consequently continuous,
even across ties and the $S$ seam; $\mathcal J$ is lower semicontinuous.
The squared distance in the statement is continuous. Its closed subset
at distance at least $\varepsilon$ has strictly positive minimum surplus,
by the equality classification. This proves the uniform assertion.
\end{proof}

\begin{lemma}[A marked coordinate]\label{lem:marked}
For each $\beta>0$, kernels with
$\max_i|\widehat u(\{i\})|\ge\beta$ satisfy
$d_a(u)\le C_\beta\Delta_r(u)$.
\end{lemma}
\begin{proof}
Split an interior field $g$ at the marked coordinate. The exact identity
$D(g)=(D(g_+)+D(g_-))/2+\E c(g_+,g_-)$ and the face production
bound give $D(g)-F(\cI(g))\ge\mathcal J$.
Also $\E c\le D(g)$. On a compact positive information interval,
small production surplus and a singleton lower bound therefore imply
arbitrarily small coordinate distance, by Lemma~\ref{lem:joining}.
Apply \eqref{eq:budget} on $[0,t]$. Some time in its middle third
has $R(s)\le3e^{Kt_{\max}}\Delta/t_{\min}$, where
$t_{\min}=-\log r_1$, $t_{\max}=-\log r_0$.
For small $\Delta$, $y(s)\in[\psi(r_0a_0)/2,\psi(a_1)]$,
and the marked singleton is at least $r_0\beta$.
The preceding compactness lemma puts $P_su$ close to the coordinate
of amplitude $\psi^{-1}(y(s))$. Equation~\eqref{eq:budget} and
Lipschitz continuity of this inverse on the positive compact interval
put that amplitude within $O(\Delta)$ of $ae^{-s}$.
Choose the distance tolerance first to enter the neighborhood in
Lemma~\ref{lem:local} via \eqref{eq:return}, then choose the production
and deficit thresholds. The local lemma gives the desired linear estimate.
Above the chosen deficit threshold the bound $d_a\le4$ suffices.
\end{proof}

\subsection{Coordinate sorting}
Sorting the two values on each coordinate fiber into increasing order
preserves the distribution of $u$, hence its mean and information.
It preserves every other singleton, and it preserves monotonicity in
coordinates already sorted. After one sweep the resulting field $g$
is increasing. Sorting does not decrease $\cI(T_ru)$: write the fiber
as $b+X_id$, and after noise in the other coordinates write
$B=T_rb$, $D=T_rd$, $D_+=T_r|d|$. Then $D_+\ge|D|$ and
\[
 f_B(q)=\tfrac12\{\psi(B+rq)+\psi(B-rq)\},\qquad f_B''(q)\ge r^2.
\]
The gain is at least $r^2\E(D_+^2-D^2)/2$.

Let $M_2=\E u^2$ and choose a coordinate with maximal final singleton
$c=\widehat g(\{i\})$. At its sorting step the initial singleton is
$b_i=\widehat u(\{i\})$ and $\E|d|=c$, because all other sorting
steps preserve that singleton. Noise spectral bounds give
\[
 \E(D_+^2-D^2)\ge(1-r^2)(c^2-b_i^2)-r^2(M_2-c^2).
\]
If $\ell=\cI(T_rg)-\cI(T_ru)$, this step's gain is at most $\ell$.
Also $d_a(g)=M_2+a^2-2ac=(M_2-c^2)+(a-c)^2$.
When $d_a(g)\le a^2/4$, $c\ge a/2$, and
\[
 d_a(u)-d_a(g)\le2a(c-|b_i|)\le4(c^2-b_i^2).
\]
Here $|b_i|\le c$, and use of a nonmaximal original singleton only
weakens the first bound. Otherwise $d_a(u)\le4\le16d_a(g)/a_0^2$.
It follows that
\begin{equation}\label{eq:compression}
 d_a(u)\le C\{d_a(g)+\ell\},\qquad
 \Delta_r(u)=\Delta_r(g)+\ell.
\end{equation}
Thus it remains to exclude increasing near-extremizers with all
singletons small. The next section supplies that exclusion and completes
Theorem~\ref{thm:compact}.
